% !TeX program = pdflatex
% !TeX TS-program = pdflatex
\documentclass[11pt,a4paper]{article}

\usepackage[T1]{fontenc}
\usepackage[utf8]{inputenc}
\usepackage{amsmath,amssymb,amsthm}
\usepackage{mathtools}
\usepackage[margin=3.2cm]{geometry}
\usepackage{microtype}
\usepackage{hyperref}
\usepackage{booktabs}
\usepackage{enumitem}
\usepackage{fancyvrb}
% Lean sources are shown through fancyvrb, not listings: listings mis-orders
% multi-byte characters adjacent to ASCII punctuation, emitting "∀(" for the
% input "(∀". Verbatim fidelity matters more here than keyword colouring.
\fvset{fontsize=\small, xleftmargin=1.2em}

% The Lean notation, for pdflatex. Declared after inputenc so the mappings also
% fire inside fancyvrb, which is where the Lean sources are set.
\DeclareUnicodeCharacter{2200}{\ensuremath{\forall}}
\DeclareUnicodeCharacter{2203}{\ensuremath{\exists}}
\DeclareUnicodeCharacter{2208}{\ensuremath{\in}}
\DeclareUnicodeCharacter{2192}{\ensuremath{\rightarrow}}
\DeclareUnicodeCharacter{2264}{\ensuremath{\leq}}
\DeclareUnicodeCharacter{2260}{\ensuremath{\neq}}
\DeclareUnicodeCharacter{2227}{\ensuremath{\wedge}}
\DeclareUnicodeCharacter{2211}{\ensuremath{\textstyle\sum}}
\DeclareUnicodeCharacter{2115}{\ensuremath{\mathbb{N}}}
\DeclareUnicodeCharacter{211D}{\ensuremath{\mathbb{R}}}
\usepackage{xcolor}

\hypersetup{colorlinks=true, linkcolor=blue!50!black, citecolor=blue!50!black,
            urlcolor=blue!50!black}


\theoremstyle{plain}
\newtheorem{theorem}{Theorem}[section]
\newtheorem{proposition}[theorem]{Proposition}
\newtheorem{lemma}[theorem]{Lemma}
\newtheorem{corollary}[theorem]{Corollary}
\newtheorem{conjecture}[theorem]{Conjecture}
\theoremstyle{definition}
\newtheorem{definition}[theorem]{Definition}
\newtheorem{remark}[theorem]{Remark}

\newcommand{\R}{\mathbb{R}}
\newcommand{\E}{\mathbb{E}}
\newcommand{\card}[1]{\lvert #1 \rvert}
\newcommand{\norm}[1]{\lVert #1 \rVert}
\newcommand{\inner}[2]{\langle #1, #2 \rangle}

\title{Exact contact numbers for packings of six to nine\\
       congruent balls, formally verified}

\author{Marek Wrzos\\[2pt]
  \normalsize\texttt{marek.wrzos.sgh@gmail.com}\\[1pt]
  \normalsize Independent researcher\\[1pt]
  \normalsize ORCID \href{https://orcid.org/0009-0005-1521-2633}{0009-0005-1521-2633}}

\date{\today}

\begin{document}
\maketitle

\begin{abstract}
The \emph{contact number} $c(n,3)$ of a packing of $n$ congruent balls in Euclidean
$3$-space is the maximum number of touching pairs the packing can have. Exact values
have been known only for $n \le 5$. Bezdek and Khan~\cite{BezdekKhan2018} record
$c(n,3) = 3n-6$ for $6 \le n \le 9$ as a conjecture, and derive it as a proposition
\emph{conditional} on the completeness of the numerical enumeration of
Arkus, Manoharan and Brenner~\cite{ArkusSIAM2011}; they state explicitly that
``the contact number $c(n,3)$ is still unknown for $n \ge 6$''.

We prove $c(6,3) = 12$, $c(7,3) = 15$, $c(8,3) = 18$ and $c(9,3) = 21$
unconditionally, and we prove in addition that every packing attaining these maxima is
minimally rigid. Since Conjecture 5.2 asserts a numerical value only for
$n = 6,\dots,9$, its numerical half is thereby settled completely, and its structural
half is settled at every $n$ for which it makes a numerical claim.
The proofs are formalized in the Lean~4 proof assistant~\cite{Lean4} on top of
\texttt{mathlib}~\cite{Mathlib2020} and are checked by Lean's kernel; they depend on
no axioms beyond Lean's three standard ones, and in particular use no compiled-code
evaluation. The geometric content is a short list of explicit lemmas about the
configuration of common neighbours of a bonded pair; the combinatorial case analysis
is delegated to kernel-verified certificates.
\end{abstract}

\section{Introduction}
\label{sec:intro}

Let $\mathcal{P}$ be a packing of $n$ congruent balls in $\E^3$: finitely many balls of
equal radius with pairwise disjoint interiors. Its \emph{contact graph} has one vertex
per ball and one edge per touching pair. The \emph{contact number} $c(n,3)$ is the
largest number of edges such a graph can have. The problem is a finite,
$3$-dimensional relative of Newton's kissing-number problem and, in the language of
distance geometry, asks for the maximum number of times the minimum distance can occur
among $n$ points of $\R^3$.

The two-dimensional analogue was settled by Harborth~\cite{Harborth1974}, who proved
that $n$ unit circles realise at most $\lfloor 3n - \sqrt{12n-3}\,\rfloor$ touching
pairs. In three dimensions no exact formula is expected. Bezdek~\cite{Bezdek2012}
proved
\[
  c(n,3) < 6n - 0.926\, n^{2/3} \qquad (n \ge 2),
\]
later improved in~\cite{BezdekReid2013}; for packings of balls of \emph{differing}
radii the corresponding average-degree bound is due to Glazyrin~\cite{Glazyrin2020}. The
leading constant $6$ reflects the kissing
number of $\R^3$ and the $n^{2/3}$ correction the surface of the cluster. For small
$n$, however, the truth is much smaller than $6n$: a contact graph that is planar has
at most $3n-6$ edges, and $3n-6$ is also the count in Maxwell's
criterion~\cite{Maxwell1864} for a minimally rigid bar-and-joint framework in space,
which is why the value $3n-6$ recurs throughout the physics literature on colloidal
clusters.

\subsection*{What was known}

Exact values of $c(n,3)$ are classical only for $n \le 5$, where they are $0,1,3,6,9$
and the arguments are elementary (for $n=5$ one needs that five points of $\R^3$ cannot
be pairwise at distance one; see Lemma~\ref{lem:k5}).

For $n \ge 6$ the literature is computational. Arkus, Manoharan and
Brenner~\cite{ArkusPRL2009,ArkusSIAM2011} enumerated packings satisfying minimal
rigidity constraints for $n \le 10$, solving the resulting distance-geometry systems
analytically for $n \le 9$ and numerically for $n = 10$. Hoy, Harwayne-Gidansky and
O'Hern~\cite{Hoy2012} extended the enumeration to $n = 11$, and
Holmes-Cerfon~\cite{HolmesCerfon2016} enumerated rigid clusters for $n \le 14$ and
maximal-contact clusters for $n \le 19$. These enumerations agree that the maximum is
$3n-6$ for $n \le 9$ and that it first exceeds $3n-6$ at $n = 10$, where the maximum is
$25 > 24$.

None of these is a proof. The authors of each are explicit about this, and the survey
of Bezdek and Khan~\cite{BezdekKhan2018} is explicit about it too: the elimination
rules of~\cite{ArkusSIAM2011} ``are susceptible to round off errors, [so] there is a
possibility that some adjacency matrices are incorrectly eliminated''; the method
of~\cite{Hoy2012} ``employ[s] Newton's method, which cannot guarantee to obtain a
solution, whenever one exists'' and rests on a Hamiltonian-path assumption shown false
by Connelly and the Demaines. The survey draws the conclusion in one sentence:

\begin{quote}
``As a result, the contact number $c(n,3)$ is still unknown for $n \ge
6$.''~\cite[\S5.1]{BezdekKhan2018}
\end{quote}

\noindent That sentence carries a footnote, which we quote in full because it is the
strongest statement in the literature against the values being open:

\begin{quote}
``According to [4], for $n \le 7$, it is possible to solve the system (3) using standard
algebraic geometry methods for all $A \in \mathcal{A}$ without filtering by geometric
rules. Arkus et al.\ [4] attempted this using the package \textsc{Singular}. Therefore,
most likely for $n = 6, 7$, the maximal contact graphs as obtained in [4] are optimal for
minimally rigid sphere packings.''~\cite[\S5.1, footnote 2]{BezdekKhan2018}
\end{quote}

\noindent Three things should be said about it. It is hedged (``most likely''). It
reaches only $n = 6, 7$, so it bears on neither $n = 8$ nor $n = 9$. And it concludes
optimality \emph{among minimally rigid packings} --- which presupposes that a
maximum-contact packing is minimally rigid, the very intuition the same survey records as
unproved: ``Up to our knowledge, there exists no proof of or counterexample to this
intuition''~\cite[\S5.1]{BezdekKhan2018}. Removing that presupposition is part of what is
done here: minimal rigidity is not assumed at any point below, it is derived
(\S\ref{sec:minrigid}). What the footnote does establish is that for $n = 6, 7$ the
algebraic step was carried out without the error-prone geometric filtering, so the
residual doubt there is narrower than for $n = 8, 9$.

What the survey \emph{does} prove is a conditional statement, which is the correct
point of comparison for the present paper:

\begin{proposition}[{\cite[Prop.~5.1]{BezdekKhan2018}}]
\label{prop:bk}
Assume that all maximal-contact minimally rigid packings of $n \le 9$ spheres are
listed in~\cite{ArkusPRL2009,ArkusSIAM2011}. Then $c(n,3) = 3n-6$ for
$n = 4,\dots,9$, and there is a minimally rigid cluster attaining it.
\end{proposition}

\begin{conjecture}[{\cite[Conj.~5.2]{BezdekKhan2018}}]
\label{conj:bk}
For $n \ge 6$, every contact graph of a packing of $n$ spheres with $c(n,3)$ contacts
is minimally rigid. Moreover $c(n,3) = 3n-6$ for $n = 6,\dots,9$.
\end{conjecture}

\subsection*{Results}

\begin{theorem}
\label{thm:main}
$c(6,3) = 12$, \quad $c(7,3) = 15$, \quad $c(8,3) = 18$, \quad $c(9,3) = 21$.
\end{theorem}

The hypothesis of Proposition~\ref{prop:bk} is removed for these $n$: nothing here
depends on any enumeration being complete, on floating-point arithmetic, or on any
rigidity assumption. The statements quantify over \emph{all} finite subsets of $\R^3$
with pairwise distances at least one --- no genericity, no boundedness, no
nondegeneracy.

Conjecture~\ref{conj:bk} has two halves with different quantifier ranges: a
\emph{structural} claim for all $n \ge 6$, and a \emph{numerical} claim for
$n = 6,\dots,9$. Theorem~\ref{thm:main} settles the numerical half for every $n$ it
concerns, and the structural half follows at those $n$ as well.

Recall the notion, which is combinatorial rather than kinematic
(\cite[Def.~1]{BezdekKhan2018}, after~\cite{ArkusSIAM2011}): a cluster of $n \ge 4$
spheres is \emph{minimally rigid} if every sphere touches at least three others and the
cluster has at least $3n-6$ contacts. The survey stresses that this is ``neither
sufficient nor necessary for rigidity'' in the sense of continuous deformability.

\begin{theorem}
\label{thm:conj52}
Conjecture~\ref{conj:bk} holds for $n = 6, 7, 8, 9$: for these $n$, $c(n,3) = 3n-6$,
the value is attained, and \emph{every} packing of $n$ congruent balls with $3n-6$
contacts is minimally rigid.
\end{theorem}

What remains of Conjecture~\ref{conj:bk} is exactly its structural clause for
$n \ge 10$, where no value $c(n,3)$ is known or conjectured by it. Theorem~\ref{thm:conj52}
still says nothing about uniqueness of the maximisers; see \S\ref{sec:notproved}.

\subsection*{Method, and what ``formally verified'' means here}

The proofs are formalized in Lean~4 and machine-checked. The argument for $n=8$
ranges over a search tree with $1\,267\,025$ nodes, each closed by one of ten geometric
obstructions; the tree is a term in the ambient logic and its verification is a kernel
reduction.

Every theorem below reports exactly \texttt{[propext, Classical.choice, Quot.sound]}.
There are no \texttt{sorry}s and no user-declared axioms. There is no use of
\texttt{native\_decide}, which would delegate evaluation to compiled code and appear as
the extra axiom \texttt{Lean.ofReduceBool}; all finite checks are performed by the Lean
kernel itself (\texttt{decide +kernel}). The geometry is the twelve lemmas of
\S\ref{sec:arsenal}, each with a short conventional proof; only the case coverage is
mechanical. Reproduction is described in \S\ref{sec:repro}.

\subsection*{Related formal work}

Computer-assisted proof in sphere packing has an established formal precedent: the
Flyspeck project verified Hales's proof of the Kepler conjecture in HOL~Light and
Isabelle~\cite{Flyspeck2017}. That concerns packing \emph{density} for infinite
packings, a different problem, but it settled the methodological question of whether a
proof of this shape can be put on a formal footing.

Closer to the present work, Hariharan, Viazovska and collaborators have formalized
Viazovska's solution of the sphere packing problem in dimension~8 in Lean, a project
begun in 2024 and completed in February 2026, with dimension~24 alongside
it~\cite{SpherePacking8Lean2026}. That is the same prover and the same corner of discrete
geometry, and it is a far larger development than this one; it is again about density
rather than contact numbers, and about a single optimal configuration in a high dimension
rather than exact extremal values for small $n$ in $\R^3$. We know of no formalization of
a contact number, or of a kissing number, in any proof assistant.

The interval branch-and-prune of \S\ref{sec:nine} has direct antecedents in that work:
Solovyev and Hales~\cite{SolovyevHales2013} verified the roughly one thousand nonlinear
inequalities of Flyspeck in HOL~Light using interval arithmetic with Taylor
approximations. The method used here is the same idea specialised and simplified --- the
inequalities are polynomial with rational coefficients, so the arithmetic can be exact
and the soundness argument purely algebraic, with no floating point and no analysis.

\subsection*{A previous claim}

A preprint of Reid~\cite{Reid2016} claims $C(6)=12$, $C(7)=15$, $C(8)=18$ and announces
values up to $C(13)=36$. It is unpublished, has one version, and is not a proof; the
argument for $n=6$ has a demonstrable gap and, more decisively, uses no geometric input
at all, which no correct argument can do. Details are in
Appendix~\ref{app:reid}.

\section{Setting}
\label{sec:setting}

We work with balls of radius $\tfrac12$, so that touching means centres at distance
exactly $1$; this differs from the normalisation of~\cite{BezdekKhan2018} (unit balls,
centres at distance $2$) by a scaling.

\begin{definition}
A finite set $X \subseteq \R^3$ is \emph{hard-core} if $\operatorname{dist}(z,w) \ge 1$
for all distinct $z,w \in X$. For $z \in X$ put
$N_X(z) = \{ w \in X : \operatorname{dist}(z,w) = 1 \}$, and
\[
  \operatorname{cc}(X) \;=\; \sum_{z \in X} \card{N_X(z)},
  \qquad
  \operatorname{E}(X) \;=\; -\tfrac12 \operatorname{cc}(X).
\]
\end{definition}

Thus $\operatorname{cc}(X)$ counts ordered touching pairs and is twice the number of
contacts, and $-\operatorname{E}(X)$ is the contact number of the packing with centre
set $X$. In this notation
\[
  c(n,3) \;=\; \max\{\, -\operatorname{E}(X) \;:\; X \text{ hard-core},\ \card X = n \,\}.
\]
The maximum is attained: the constructions of \S\ref{sec:constructions} realise it, so
no compactness argument is needed.

In Lean these are, verbatim:

\begin{Verbatim}
abbrev E3 := EuclideanSpace ℝ (Fin 3)

def HardCore (X : Finset E3) : Prop :=
  ∀ z ∈ X, ∀ w ∈ X, z ≠ w → 1 ≤ dist z w

noncomputable def neighbors (X : Finset E3) (z : E3) : Finset E3 :=
  X.filter (fun w => dist z w = 1)

noncomputable def contactCount (X : Finset E3) : ℕ :=
  ∑ z ∈ X, (neighbors X z).card

noncomputable def energy (X : Finset E3) : ℝ := -(contactCount X : ℝ) / 2
\end{Verbatim}

\noindent On top of these, \texttt{Interface.lean} states the results in the language of
the contact-number problem, and it is those statements that
Theorems~\ref{thm:main} and~\ref{thm:conj52} refer to and that the
trust-base check reports:

\begin{Verbatim}
noncomputable def contacts (X : Finset E3) : ℕ := contactCount X / 2

def realised (n : ℕ) : Set ℕ :=
  {k | ∃ X : Finset E3, HardCore X ∧ X.card = n ∧ contacts X = k}

def MinimallyRigid (X : Finset E3) : Prop :=
  (∀ v ∈ X, 3 ≤ (neighbors X v).card) ∧ 3 * X.card - 6 ≤ contacts X

theorem contactNumber_six   : IsGreatest (realised 6) 12
theorem contactNumber_seven : IsGreatest (realised 7) 15
theorem contactNumber_eight : IsGreatest (realised 8) 18
theorem contactNumber_nine  : IsGreatest (realised 9) 21

theorem conjecture52_nine :
    IsGreatest (realised 9) 21 ∧
    ∀ X : Finset E3, HardCore X → X.card = 9 → contacts X = 21 →
      MinimallyRigid X
\end{Verbatim}

\noindent \texttt{IsGreatest S c} is \texttt{mathlib}'s ``$c$ is the greatest element of
$S$'', so \texttt{contactNumber\_nine} is literally $c(9,3) = 21$, and
\texttt{MinimallyRigid} is Definition~1 of~\cite{BezdekKhan2018} written out. Each of
these reports the three standard axioms and nothing more.

We note where the content of \texttt{conjecture52\_nine} lies. Of the two clauses of
minimal rigidity, the second --- at least $3n-6$ contacts --- is immediate once the
packing is assumed to attain $c(n,3) = 3n-6$, and is stated only because Definition~1
includes it. Everything to be proved is in the first clause, that no ball touches fewer
than three others; that is \S\ref{sec:minrigid}.

\begin{remark}
Two small facts are used silently throughout. First, $\operatorname{cc}(X)$ is even
(the handshake lemma: bonds come in ordered pairs), which lets us step from an odd
lower bound to the next even one. Second, deleting a vertex $v$ removes exactly
$2\card{N_X(v)}$ from $\operatorname{cc}$, which is what makes the recursion of
\S\ref{sec:eight} work.
\end{remark}

\section{The local arsenal}
\label{sec:arsenal}

All geometry used in this paper is in the following twelve statements. Every one is
a statement about a bounded number of points of $\R^3$, and every one is proved in the
formalization by an explicit algebraic identity --- there is no appeal to spherical
geometry, convexity, or topology anywhere in the development.

\subsection{Five points}

\begin{lemma}[\texttt{pairwise\_unit\_card\_le\_four}]
\label{lem:k5}
If $S \subseteq \R^3$ is finite and $\operatorname{dist}(u,v) = 1$ for all distinct
$u,v \in S$, then $\card S \le 4$.
\end{lemma}

\begin{proof}
Suppose $p_0,\dots,p_4$ are pairwise at distance $1$ and set $q_i = p_i - p_0$ for
$i = 1,\dots,4$. Then $\norm{q_i} = 1$ and, expanding
$\norm{q_i - q_j}^2 = 1$, $\inner{q_i}{q_j} = \tfrac12$ for $i \ne j$. Four vectors in
$\R^3$ are linearly dependent, so $\sum_i g_i q_i = 0$ for some $g \ne 0$, whence
\[
  0 \;=\; \Bigl\lVert \sum_i g_i q_i \Bigr\rVert^2
    \;=\; \sum_i g_i^2 + \tfrac12\sum_{i \ne j} g_i g_j
    \;=\; \tfrac12 \sum_i g_i^2 + \tfrac12 \Bigl(\sum_i g_i\Bigr)^2 \;>\; 0,
\]
a contradiction.
\end{proof}

This alone gives $c(4,3)=6$ and $c(5,3)=9$.

\subsection{The ring of a bonded pair}

The single geometric idea behind $n = 6,7,8$ is that the common neighbours of a bonded
pair are severely constrained. Let $v,w$ be at distance $1$ and let $m = \tfrac12(v+w)$
be the midpoint. If $u$ is at distance $1$ from both, then $z = u - m$ satisfies
\[
  \inner{z}{w-v} = 0,
  \qquad
  \norm{z}^2 = \tfrac34 .
\]
So all common neighbours, recentred at $m$, lie on a circle of radius $\sqrt3/2$ in the
plane $(w-v)^\perp$. We call this the \emph{ring} of the bond $vw$. For two ring
vectors $z,z'$,
\[
  \norm{z - z'}^2 = \tfrac32 - 2\inner{z}{z'},
\]
so the hard-core condition $\norm{z-z'} \ge 1$ is exactly $\inner{z}{z'} \le \tfrac14$,
and $z,z'$ are themselves in contact exactly when $\inner{z}{z'} = \tfrac14$. The ring
is therefore a set of vectors of squared norm $3/4$ in a $2$-plane, pairwise at inner
product at most $1/4$, and its contact graph is the graph of \emph{equality}.

\begin{lemma}[\texttt{ring\_card\_le\_five}, \texttt{common\_neighbors\_le\_five}]
\label{lem:ringfive}
A ring has at most five elements. Consequently, two bonded particles of a hard-core
configuration have at most five common neighbours.
\end{lemma}

\begin{proof}
Writing ring vectors in polar form, the cosine of the angle between two of them is
$\inner{z}{z'} / (3/4) \le 1/3$, so consecutive azimuths differ by more than
$\arccos(1/3) > 70^\circ$. Six such directions would need more than $423^\circ$.
\end{proof}

Both bounds are sharp: face-centred and hexagonal close packings realise four common
neighbours, the pentagonal bipyramid five.

\begin{lemma}[\texttt{ring\_no\_triangle}]
\label{lem:notri}
Three ring vectors pairwise at inner product exactly $1/4$ do not exist.
\end{lemma}

\begin{proof}
Their Gram matrix has all diagonal entries $3/4$ and all off-diagonal entries $1/4$,
hence determinant $5/16 \ne 0$; but they lie in a $2$-plane, so the Gram matrix has
rank at most two.
\end{proof}

\begin{lemma}[\texttt{ring\_no\_square}]
\label{lem:nosq}
Four ring vectors $a,b,c,d$ with
$\inner{a}{c} = \inner{a}{d} = \inner{b}{c} = \inner{b}{d} = 1/4$
and $\inner{a}{b} \le 1/4$, $\inner{c}{d} \le 1/4$ do not exist.
\end{lemma}

\begin{proof}
Both $c$ and $d$ are orthogonal to $a - b$, which is a nonzero vector of the $2$-plane,
so $c$ and $d$ are parallel; as both have squared norm $3/4$,
$\inner{c}{d} = \pm 3/4$. The value $+3/4$ forces $c = d$, and $-3/4$ violates
$\inner{c}{d} \le 1/4$.
\end{proof}

\begin{lemma}[\texttt{ring\_two\_positions}]
\label{lem:twopos}
Let $c,x,y$ be ring vectors with $\inner{c}{x} = \inner{c}{y} = \alpha$. Then
\[
  \inner{x}{y} = \tfrac34
  \qquad\text{or}\qquad
  \inner{x}{y} = \tfrac{8}{3}\alpha^2 - \tfrac34 .
\]
In particular $\alpha = \tfrac14$ gives $\inner{x}{y} \in \{\tfrac34, -\tfrac{7}{12}\}$,
and $\alpha = -\tfrac{7}{12}$ gives $\inner{x}{y} \in \{\tfrac34, \tfrac{17}{108}\}$.
\end{lemma}

\begin{proof}
Choose coordinates in the $2$-plane with $c = (\tfrac{\sqrt3}{2},0)$ and write
$x = \tfrac{\sqrt3}{2}(\cos\theta,\sin\theta)$. Then $\inner{c}{x} = \tfrac34\cos\theta$,
so $\cos\theta = \tfrac{4\alpha}{3}$ determines $\theta$ up to sign: $x$ and $y$ are
either equal (giving $\inner{x}{y} = \norm{x}^2 = \tfrac34$) or reflections of each
other in the line $\R c$, in which case
$\inner{x}{y} = \tfrac34\cos 2\theta = \tfrac34(2\cos^2\theta - 1)
= \tfrac{8}{3}\alpha^2 - \tfrac34$.
\end{proof}

\begin{lemma}[\texttt{ring\_degree\_le\_two}]
\label{lem:ringdeg}
A ring vector is in contact with at most two other ring vectors. Hence the contact
graph of a ring is a disjoint union of paths.
\end{lemma}

\begin{proof}
Suppose $x,y,z$ are ring vectors with $\inner{c}{x} = \inner{c}{y} = \inner{c}{z} =
1/4$. By Lemma~\ref{lem:twopos} each pair among $x,y,z$ has inner product $3/4$ or
$-7/12$; the value $3/4$ means coincidence, so all three pairs are at $-7/12$. Applying
Lemma~\ref{lem:twopos} again at $\alpha = -7/12$ puts any two of them at $3/4$ or
$17/108$, and neither is $-7/12$.
\end{proof}

\begin{lemma}[\texttt{ring\_no\_pentagon}]
\label{lem:nopent}
A five-cycle of contacts among ring vectors, with the two pairs skew at $z_1$
hard-core, does not exist.
\end{lemma}

\begin{proof}
Chaining Lemma~\ref{lem:twopos} around the cycle forces
$\inner{z_1}{z_3} = \inner{z_1}{z_4} = -7/12$, and then $z_3, z_4$ must sit at inner
product $3/4$ or $17/108$; but they are adjacent in the cycle, so at $1/4$.
\end{proof}

\subsection{Common neighbours of a triple}

\begin{lemma}[\texttt{common\_neighbors\_triple\_le\_two}]
\label{lem:triple}
Three distinct points of $\R^3$ have at most two common neighbours. No hard-core
hypothesis is needed.
\end{lemma}

\begin{proof}
Let $a,b,c$ be distinct and $w_1,w_2,w_3$ at distance $1$ from each. Differences
$w_i - w_j$ are orthogonal to $b-a$ and $c-a$. The four vectors
$b-a$, $c-a$, $w_2-w_1$, $w_3-w_1$ lie in $\R^3$ and so admit a nontrivial relation,
which splits into two vanishing halves by orthogonality; the Gram kernel of either half
collapses two of the five points, contradicting distinctness.
\end{proof}

\subsection{Shell obstructions}

For $n = 8$ three further patterns are needed, each ruling out a specific contact
pattern around a particle of high degree. Throughout, $ij$ abbreviates the pair
$\{u_i, u_j\}$.

If a particle $c$ is in contact with $u_0,\dots,u_5$, then recentring at $c$ makes the
six differences $u_i - c$ unit vectors, and two of them are in contact exactly when
their inner product is $1/2$. Six unit vectors span at most three dimensions, so their
Gram matrix is singular in a strong sense:

\begin{lemma}[\texttt{gram\_det\_zero}]
\label{lem:gram}
Any four vectors of $\R^3$ have singular $4 \times 4$ Gram matrix.
\end{lemma}

\begin{lemma}[shell obstructions]
\label{lem:patterns}
Let $X$ be hard-core and let $c, u_0, \dots, u_5 \in X$ be distinct. None of the
following configurations occurs.
\begin{description}[leftmargin=0pt, itemsep=3pt, font=\normalfont\itshape]
\item[$P_1$.] $c$ is in contact with every $u_i$, and the contacts among the $u_i$
  include $03$, $04$, $05$, $12$, $14$, $15$, $23$, $24$, $34$.
\item[$P_4$.] $c$ is in contact with every $u_i$, and the contacts among the $u_i$
  include $01$, $04$, $05$, $12$, $13$, $23$, $25$, $34$, $45$.
\item[$P_5$.] $c$ is in contact with $u_1, \dots, u_5$, not necessarily with $u_0$, and
  the contacts among the $u_i$ include $03$, $04$, $05$, $12$, $14$, $15$, $23$, $25$, $34$.
\end{description}
\end{lemma}

\begin{proof}[Proof sketch]
For $P_1$ and $P_4$ the recentred shell spans at most three dimensions, so by
Lemma~\ref{lem:gram} every $4 \times 4$ Gram determinant among the six shell vectors
vanishes. The prescribed contacts fix most entries of the Gram matrix, and the
vanishing determinants become polynomial equations in the few remaining ones. For
$P_1$ they force $x_{02} = x_{13} = -1/3$ and $x_{01} = -7/18$, after which
$\det G[0,1,4,5]$ has roots $1$ and $7/11$, both above the hard-core ceiling $1/2$.
For $P_4$ the free entries lie on a $6$-cycle; the determinants of opposite pairs
factor as $(a-1)(b-1)(ab+a+b)$, so hard-core forces $ab + a + b = 0$ exactly, and the
adjacent-pair determinants $12a^2 - 8ab - 4a + 12b^2 - 4b - 5 = 0$ then admit no
admissible solution.

$P_5$ has no vertex adjacent to all six, so there is no common centre to recentre at
and the argument above does not apply directly. It is run instead on the differences
$u_i - u_j$, whose Gram matrix is again singular by Lemma~\ref{lem:gram}, with the
contact conditions expressed in those coordinates.

The three arguments are \texttt{pattern\_p1\_impossible},
\texttt{pattern\_p4\_impossible} and \texttt{pattern\_p5\_impossible} in the
formalization, which carries them out in full.
\end{proof}

\section{Six particles}
\label{sec:six}

This case needs no enumeration and is the model for the rest.

\begin{theorem}
\label{thm:six}
$c(6,3) = 12$.
\end{theorem}

\begin{proof}
For the upper bound, suppose $X$ is hard-core with $\card X = 6$ and
$\operatorname{cc}(X) \ge 25$. Since $\operatorname{cc}$ is even, $\operatorname{cc}(X)
\ge 26$. Every degree is at most $5$, so if at most one vertex had degree $5$ the sum
would be at most $5 + 5\cdot 4 = 25$; hence there are two vertices $e \ne f$ of degree
$5$.

Being of degree $5$ in a $6$-element set, $e$ and $f$ are each bonded to everything
else --- in particular to each other --- and the remaining four particles are common
neighbours of the bond $ef$. Nine of the thirteen edges are thus accounted for
($ef$, and four from each of $e$ and $f$), so at least four edges join the four
ring vertices. A graph on four vertices with four edges contains a cycle, necessarily
a triangle or a four-cycle, and these are excluded by Lemmas~\ref{lem:notri}
and~\ref{lem:nosq} respectively (the diagonals of the four-cycle are hard-core, as
required).

For the lower bound, the regular octahedron with vertices
$(\pm\tfrac{\sqrt2}{2},0,0)$, $(0,\pm\tfrac{\sqrt2}{2},0)$,
$(0,0,\pm\tfrac{\sqrt2}{2})$ is hard-core --- adjacent vertices are at distance $1$ and
antipodal ones at $\sqrt2$ --- and has $12$ contacts.
\end{proof}

\section{Seven particles}
\label{sec:seven}

\begin{theorem}
\label{thm:seven}
$c(7,3) = 15$.
\end{theorem}

Here counting no longer suffices and a case analysis begins. Suppose $X$ is hard-core
with $\card X = 7$ and at least $16$ contacts. Its contact graph is $K_7$ minus at most
five of the $21$ pairs. Enumerating the $\binom{21}{5} = 20\,349$ five-element sets $R$
of pairs, it is enough to show that no hard-core realisation exists whose non-contacts
are contained in $R$ --- taking a \emph{superset} of the actual non-contacts is what
makes the enumeration finite and every obstruction monotone in the edge set.

For each $R$ we exhibit one witness of one of six kinds, each an instance of a lemma of
\S\ref{sec:arsenal}:

\begin{center}
\begin{tabular}{clc}
\toprule
kind & obstruction & lemma \\
\midrule
0 & five pairwise bonded vertices & \ref{lem:k5} \\
1 & a bond, a common neighbour bonded to three others & \ref{lem:ringdeg} \\
2 & a bond with three pairwise bonded common neighbours & \ref{lem:notri} \\
3 & a bond with a four-cycle among its common neighbours & \ref{lem:nosq} \\
4 & a bond with a five-cycle among its common neighbours & \ref{lem:nopent} \\
5 & three vertices with three common neighbours & \ref{lem:triple} \\
\bottomrule
\end{tabular}
\end{center}

A witness is a single natural number packing the kind and up to seven vertex labels in
base $8$; a Boolean function \texttt{killcheck} verifies that the labelled vertices are
pairwise adjacent as the kind requires, given $R$. The list of $20\,349$ witnesses is
machine-generated; that \emph{every} case carries a valid one is the single Boolean
identity

\begin{Verbatim}
theorem all_killed : killAll combos emin7Certs = true
\end{Verbatim}

\noindent verified by the kernel in seven chunks of $3000$ (a monolithic fold exceeds
the kernel's recursion depth) and reassembled. A separate soundness lemma,
\texttt{killcheck\_sound}, states that a verified witness contradicts the existence of
a hard-core realisation; it is proved once per kind by discharging to the corresponding
lemma of \S\ref{sec:arsenal}. Together these give $\operatorname{cc}(X) \le 30$.

The lower bound is the pentagonal bipyramid of \S\ref{sec:constructions}.

\section{Eight particles}
\label{sec:eight}

\begin{theorem}
\label{thm:eight}
$c(8,3) = 18$.
\end{theorem}

At $n=8$ the complement has up to nine pairs and $\binom{28}{9} \approx 6.9$ million
cases, so the flat enumeration of \S\ref{sec:seven} is replaced by a search tree. Two
ingredients make it small enough.

\paragraph{A recursion off $n=7$.} If $X$ is hard-core with $\card X = 8$ and
$\operatorname{cc}(X) \ge 37$, then every particle has at least four bonds: deleting a
particle of degree at most three leaves seven particles carrying at least $31$ ordered
contacts, contradicting Theorem~\ref{thm:seven}. This is
\texttt{degree\_ge\_four}, and it is why $n=8$ genuinely rests on $n=7$.

\paragraph{Orderly generation.} We may enumerate the particles in order of
non-increasing degree. A branch in which the maximal possible degree at position $i$
falls below the minimal possible degree at position $i+1$ is therefore dead, which
prunes the tree substantially and costs one extra certificate kind.

The search branches on the $28$ pairs in a fixed vertex-major order, deciding each to
be a contact or a non-contact, and closes a branch as soon as one of ten certificates
applies. Three of the six kinds of \S\ref{sec:seven} carry over --- $K_5$, a ring vertex
of degree three, and a triple with three common neighbours --- joined by an edge with
six common neighbours (Lemma~\ref{lem:ringfive}), the three shell patterns of
Lemma~\ref{lem:patterns}, a vertex with four decided non-contacts (dead by
\texttt{degree\_ge\_four}), ten decided non-contacts (dead by counting), and the
degree-order rule above. The ring four- and five-cycles of \S\ref{sec:seven} are not
needed as separate kinds here. The full list is Appendix~\ref{app:certs}, and the
resulting tree has $1\,267\,025$ nodes.

The tree is serialized as a list of naturals --- $0$ opens a branch, an odd entry is a
certificate leaf --- and a walker
\begin{Verbatim}
def walk : ℕ → List ℕ → ℕ → List ℕ → List ℕ → Option (List ℕ)
\end{Verbatim}
verifies simultaneously that the shape is a complete binary decision tree over the
pairs and that every leaf certificate holds for the edge and non-edge sets decided on
the path to it. Coverage of all contact patterns is therefore intrinsic to the shape:
there is nothing to check separately.

The kernel cannot evaluate the whole walk at once --- the monolithic reduction
overflows the kernel stack at roughly $1.27$ million nested frames --- so it is checked
in $264$ subtree-aligned chunks of at most $8000$ nodes and recomposed with three
structural lemmas (\texttt{walk\_append}, \texttt{walk\_branch},
\texttt{walk\_fuel\_mono}). Soundness is \texttt{walk\_impossible}: a successful walk of
any tree is incompatible with a hard-core eight-particle configuration of at least $38$
ordered contacts. It is proved by induction on the recursion, with one case per
certificate kind.

The lower bound is the capped pentagonal bipyramid of \S\ref{sec:constructions}.

\section{Nine particles}
\label{sec:nine}

\begin{theorem}
\label{thm:nine}
$c(9,3) = 21$.
\end{theorem}

At $n=9$ there are $\binom92 = 36$ pairs and the complement of a $22$-contact graph has
up to $14$ of them. Three things change relative to \S\ref{sec:eight}.

\paragraph{The recursion, one level up.} If $X$ is hard-core with $\card X = 9$ and
$\operatorname{cc}(X) \ge 43$, then every particle has at least four bonds: deleting one
of degree at most three leaves eight particles with at least $37$ ordered contacts,
contradicting Theorem~\ref{thm:eight}. This is \texttt{degree\_ge\_four\_9}, and it is
why $n=9$ rests on $n=8$ exactly as $n=8$ rests on $n=7$.

\paragraph{Canonical enumeration replaces degree sorting.} The degree-order rule of
\S\ref{sec:eight} is too weak here. Instead we enumerate only the lexicographically
maximal representative of each isomorphism class of decided bond patterns: a branch dies
as soon as some vertex permutation produces a strictly larger bond code. Soundness is
\texttt{exists\_lexmax\_enum} together with \texttt{bondCode\_lt}, and the certificate
witnessing a kill is the permutation. This is orderly generation rather than a partial
symmetry break, and it keeps the tree at $131\,325$ nodes, three orders of magnitude
below the raw search.

\paragraph{Twenty-three kill kinds.} Nine of the ten of \S\ref{sec:eight} survive
unchanged; the degree-order rule is replaced by the canonicity rule just described. The
thirteen new ones
close shell and near-shell patterns that first occur at $n=9$, including configurations
around vertices of degree seven and eight that the lemmas of \S\ref{sec:arsenal} do not
reach. They are certified by two mechanisms.

\begin{itemize}[leftmargin=*]
\item \emph{Positivstellensatz certificates.} For most patterns, the Gram-determinant
  system is infeasible for a reason expressible as an explicit identity: a nonnegative
  combination of the constraint polynomials and squares that sums to a negative
  constant. Verifying such a certificate is a polynomial identity check, which the
  kernel does by \texttt{ring}-style normalisation over $\mathbb{Q}$.
\item \emph{Interval branch-and-prune over $\mathbb{Q}$} (\texttt{IntervalBP}). Where no
  low-degree certificate was found, infeasibility is established by exhaustive
  subdivision of a rational box. Polynomials are carried as lists of
  (rational coefficient, exponent vector); a Taylor shift moves the origin to the corner
  of a box, after which a sign inspection of the shifted coefficients yields an exact
  rational lower bound on the box. Subdivision proceeds dyadically until every leaf box
  is excluded. Everything is algebraic: the soundness proof uses no analysis, no
  floating point, and no interval library --- only the shift identity and monotonicity
  of the bound.
\end{itemize}

One pattern (\texttt{Q28}) required its own branch-and-prune tree, split across nine
modules for the kernel.

\begin{remark}
\texttt{IntervalBP} is not specific to contact numbers: it is a kernel-checkable
framework for refuting systems of polynomial inequalities over rational boxes, and
applies wherever a semialgebraic set must be excluded that resists a
Positivstellensatz certificate of low degree.
\end{remark}

The tree is verified and recomposed exactly as in \S\ref{sec:eight}, giving
$\operatorname{cc}(X) \le 42$. The lower bound is the doubly-capped pentagonal bipyramid
of \S\ref{sec:constructions}.

\section{Minimal rigidity of the maximisers}
\label{sec:minrigid}

Theorem~\ref{thm:conj52} follows from Theorem~\ref{thm:main} by one further step, which
costs nothing beyond the cases already proved.

Let $X$ be hard-core with $\card X = n$ and exactly $3n-6$ contacts. The second clause
of minimal rigidity, ``at least $3n-6$ contacts'', holds with equality. For the first,
suppose some $v \in X$ touches at most two others. Deleting $v$ removes exactly
$2\deg(v) \le 4$ from $\operatorname{cc}$, so
\[
  \operatorname{cc}(X \setminus \{v\}) \;\ge\; 2(3n-6) - 4 \;=\; 2\bigl(3(n-1)-6\bigr) + 2 ,
\]
that is, $X \setminus \{v\}$ is a hard-core packing of $n-1$ balls with at least
$3(n-1)-6+1$ contacts, exceeding $c(n-1,3)$. So each case rests on its predecessor:
$n=9$ on $c(8,3)=18$, $n=8$ on $c(7,3)=15$, $n=7$ on $c(6,3)=12$, and $n=6$ on
$c(5,3)=9$, which is Lemma~\ref{lem:k5}. (The case $n=6$ needs no hard-core hypothesis
at all, since $c(5,3)=9$ is proved by counting.)

This is the same erase-a-vertex recursion as \texttt{degree\_ge\_four} in
\S\ref{sec:eight}, run one step lower; in the formalization it is
\texttt{minDegree\_six} through \texttt{minDegree\_nine}, and the packaged statements
\texttt{conj52\_six} through \texttt{conj52\_nine}.

\begin{remark}
The structural half of Conjecture~\ref{conj:bk} is asserted for \emph{all} $n \ge 6$,
and in that generality it is out of reach of this method --- or of any finite
computation --- because $c(n,3)$ is unknown for every $n \ge 10$. What is proved here
are its instances at $n = 6,\dots,9$, which is every $n$ for which
Conjecture~\ref{conj:bk} itself asserts a value.
\end{remark}

\section{The extremal configurations}
\label{sec:constructions}

Each of the following is verified in the formalization to be hard-core and to have the
stated number of contacts, with exact algebraic coordinates.

\begin{itemize}[leftmargin=*]
\item $n=6$: the regular \textbf{octahedron}, vertices
  $(\pm\tfrac{\sqrt2}{2},0,0)$ and cyclic; $12$ contacts.
\item $n=7$: the \textbf{pentagonal bipyramid} with unit edges. With
  $t = \sqrt{10 - 2\sqrt5}$, the equatorial pentagon is
  $\bigl(\tfrac{(5+\sqrt5)t}{20},0,0\bigr)$ and its rotations, and the apexes are
  $\bigl(0,0,\pm\tfrac{\sqrt5\,t}{10}\bigr)$; $15$ contacts. The apexes are at
  distance $\tfrac{\sqrt5\,t}{5} \approx 1.0515 > 1$ and so are \emph{not} in contact.
\item $n=8$: the \textbf{capped pentagonal bipyramid}, obtained by placing an eighth
  particle in contact with one equatorial edge and one apex; $18$ contacts.
\item $n=9$: the \textbf{doubly-capped pentagonal bipyramid}, with a ninth particle
  capping a second face; $21$ contacts.
\end{itemize}

Hard-coreness and the contact count are theorems about the exact coordinates, not
numerical checks.

\section{What is not proved}
\label{sec:notproved}

We are explicit, because Conjecture~\ref{conj:bk} contains more than
Theorem~\ref{thm:main} delivers.

\begin{itemize}[leftmargin=*]
\item \textbf{Not minimal rigidity in general.} We prove the structural half of
  Conjecture~\ref{conj:bk} at $n = 6,\dots,9$ (\S\ref{sec:minrigid}); the conjecture
  asserts it for every $n \ge 6$, and for $n \ge 10$ it remains open. This is the only
  part of Conjecture~\ref{conj:bk} that survives, and no finite computation can reach
  it, since $c(n,3)$ is unknown for every $n \ge 10$.
\item \textbf{Not rigidity.} ``Minimally rigid'' is a counting condition, and the
  survey notes it is neither sufficient nor necessary for rigidity; indeed
  \cite[Fig.~1]{BezdekKhan2018} exhibits a minimally rigid $9$-sphere packing that is
  not rigid. Nothing here asserts that maximum-contact packings are rigid.
\item \textbf{Not uniqueness.} We prove that the maximum is attained by the listed
  configurations, not that they are the only maximisers. The enumerations
  of~\cite{ArkusSIAM2011,HolmesCerfon2016} indicate several combinatorially distinct
  maximisers at $n = 8$; identifying them rigorously is a separate problem.
\item \textbf{Not $n \ge 10$.} In particular $c(10,3)$, where the pattern $3n-6$ is
  believed to break with $c(10,3) = 25$, is untouched.
\item \textbf{Not a continuous-potential statement.} The model here is exact contact.
  Transferring these bounds to narrow-well potentials requires quantitative,
  $\varepsilon$-tolerant versions of the lemmas of \S\ref{sec:arsenal}; the
  certificates carry margins that make this plausible, but it is not done.
\end{itemize}

\section{Reproducibility}
\label{sec:repro}

The development is at
\url{https://github.com/Pendzoncymisio/contact-numbers-lean} and is archived at
\href{https://doi.org/10.5281/zenodo.21935524}{\texttt{doi:10.5281/zenodo.21935524}},
which resolves to the current version. It depends only on Lean~4
(toolchain \texttt{leanprover/lean4:v4.32.0}) and \texttt{mathlib} at the matching
release; there are no other dependencies. After \texttt{lake build}, the command

\begin{Verbatim}
lake env lean Axioms.lean
\end{Verbatim}

\noindent prints the axiom dependencies of every theorem in Theorems~\ref{thm:main}
and~\ref{thm:conj52}, together with the full elaborated statements and the definitions
they use. The expected output is \texttt{[propext, Classical.choice, Quot.sound]} for
each.

The development is $95$ modules and $50\,216$ lines. The mathematics is $16$ of those
modules and $9\,352$ lines: the definitions, the lemmas of \S\ref{sec:arsenal}, the
reductions and their soundness proofs, and \texttt{IntervalBP}. The other $79$ modules,
$40\,864$ lines and $12.7$\,MB, are machine-emitted --- the certificates of
\S\ref{sec:seven}, the serialized trees of \S\S\ref{sec:eight}--\ref{sec:nine}, the
Positivstellensatz witnesses, and the branch-and-prune subdivisions, each split across
several modules so that no single one exceeds what a continuous-integration job can
build. There are $6\,437$ kernel evaluation sites (\texttt{decide +kernel}) and no other
computational steps.

The certificate generators (Python) are included for provenance. They are \emph{not}
part of the trust base: their output is data, and its correctness is established by the
kernel checks described above, not by the generators being right.

\section{Outlook}
\label{sec:outlook}

The obstruction-plus-certified-tree architecture is not specific to any one $n$. The
binding constraints in extending it are the size of the search tree and the fact that
each new $n$ introduces shell patterns not reached by the existing lemmas ---
\texttt{IntervalBP} was built precisely because that gap opened at $n=9$.

The next target is $n = 10$, and it is qualitatively different from everything above.
There the pattern $3n-6$ is believed to break, with $c(10,3) = 25 > 24$. Proving it
requires a $25$-contact construction \emph{and} a matching upper bound, and the
recursion used throughout this paper --- delete a low-degree particle and appeal to the
previous case --- no longer bites in the same way, because $25$ exceeds what
$c(9,3) = 21$ constrains. It would be the first rigorous determination of a contact
number exceeding $3n-6$, and the first point at which the small-$n$ pattern is known,
rather than believed, to fail.

Beyond the exact ladder, the asymptotic problem is a separate programme with separate
tools: $c(n,3) = 6n - \Theta(n^{2/3})$ is provable by counting and isoperimetry rather
than enumeration, and sharpening the constant in the $n^{2/3}$ term --- the discrete
Wulff-shape problem --- is where a proof of three-dimensional crystallization would
have to begin. No amount of exact small-$n$ data approaches it.

\section*{Use of AI assistance}

All mathematical content of this paper --- statements, proofs, certificate generators,
and the Lean~4 formalization --- was developed in collaboration with Claude (Anthropic),
operating as an autonomous agent under the author's direction. The author is solely
responsible for the work. Every theorem asserted here is checked by the Lean~4 kernel
and depends only on Lean's three standard axioms (\texttt{propext},
\texttt{Classical.choice}, \texttt{Quot.sound}); readers are invited to verify this
independently by running \texttt{lake env lean Axioms.lean} in the accompanying
repository, as described in \S\ref{sec:repro}. All bibliographic references have been
checked against the primary sources.

\appendix

\section{The certificate kinds}
\label{app:certs}

A certificate is a single natural number. Its kind is the residue modulo $16$; the
payload is read off as base-$8$ digits above that, so field $k$ is
$\lfloor \mathrm{cert}/(16\cdot 8^{k})\rfloor \bmod 8$, giving up to eight vertex labels
of a nine-point configuration. A Boolean function (\texttt{checkCert}) tests the payload
against the edges and non-edges decided on the path to the leaf; a soundness lemma, one
case per kind, turns a passing test into a contradiction with the lemma named in the
last column. Kinds are monotone in the edge set: they inspect only pairs already decided
to be contacts, which is what allows a branch to be closed before all pairs are decided.

\medskip
\noindent\textbf{The ten kinds for $n = 8$.}

\begin{center}
\small
\begin{tabular}{cl p{5.1cm} l}
\toprule
kind & payload & the branch is dead because & discharged by \\
\midrule
0 & $v$                       & $v$ has four decided non-contacts, so degree $\le 3$
                              & \texttt{degree\_ge\_four} \\
1 & ---                       & ten decided non-contacts, so at most $18$ contacts
                              & counting \\
2 & $a,b,c,d,e$               & five pairwise touching balls
                              & Lemma~\ref{lem:k5} \\
3 & $u,v,c,x,y,z$             & a ring vertex touching three others
                              & Lemma~\ref{lem:ringdeg} \\
4 & $u,v,w_1,\dots,w_6$       & a contact with six common neighbours
                              & Lemma~\ref{lem:ringfive} \\
5 & $a,b,c,w_1,w_2,w_3$       & three balls with three common neighbours
                              & Lemma~\ref{lem:triple} \\
6 & $u_0,\dots,u_5,c$         & shell pattern $P_1$
                              & Lemma~\ref{lem:patterns} \\
7 & $u_0,\dots,u_5,c$         & shell pattern $P_4$
                              & Lemma~\ref{lem:patterns} \\
8 & $u_0,\dots,u_5,c$         & shell pattern $P_5$
                              & Lemma~\ref{lem:patterns} \\
9 & $i$                       & degree at position $i+1$ exceeds that at position $i$
                              & \texttt{exists\_sorted\_enum} \\
\bottomrule
\end{tabular}
\end{center}

\medskip
\noindent\textbf{The twenty-three kinds for $n = 9$.} Kinds $0$--$8$ are as above. Kind
$9$ becomes the canonicity rule of \S\ref{sec:nine}, whose payload is a vertex
permutation and which is discharged by \texttt{exists\_lexmax\_enum} together with
\texttt{bondCode\_lt}. Kinds $10$--$22$ are further contact patterns, each ruled out by
its own lemma:

\begin{center}
\small
\begin{tabular}{ll@{\qquad}ll}
\toprule
kind & lemma & kind & lemma \\
\midrule
10 & \texttt{pattern\_p6\_impossible}      & 17 & \texttt{pattern\_q10\_impossible} \\
11 & \texttt{pattern\_obs7\_impossible}    & 18 & \texttt{pattern\_q1\_impossible} \\
12 & \texttt{pattern\_obs3\_impossible}    & 19 & \texttt{pattern\_q12\_impossible} \\
13 & \texttt{pattern\_obs0\_impossible}    & 20 & \texttt{pattern\_q28\_impossible} \\
14 & \texttt{pattern\_class23\_impossible} & 21 & \texttt{pattern\_q30\_impossible} \\
15 & \texttt{pattern\_q0\_impossible}      & 22 & \texttt{pattern\_q37\_impossible} \\
16 & \texttt{pattern\_q4\_impossible}      &    & \\
\bottomrule
\end{tabular}
\end{center}

\noindent Kinds $10$--$14$ are closed by rational determinant chains or by explicit
Positivstellensatz certificates. Among the \texttt{q} patterns, \texttt{q0},
\texttt{q4} and \texttt{q10} carry Positivstellensatz certificates; \texttt{q1},
\texttt{q12}, \texttt{q28}, \texttt{q30} and \texttt{q37} need the interval
branch-and-prune of \S\ref{sec:nine} and carry their own subdivision trees, that of
\texttt{q28} split across nine modules.

\section{On the preprint of Reid}
\label{app:reid}

Reid~\cite{Reid2016} claims to prove $C(6)=12$, $C(7)=15$, $C(8)=18$ and announces
proofs of $C(9)=21$ through $C(13)=36$. We record why the argument does not establish
its conclusions.

The proof of $C(6)=12$ proceeds entirely by counting degrees. Assuming a packing of six
balls with $13$ contacts, it deduces that at least two balls have degree $5$, says
``no other spheres have degree $5$, so at most $2$ spheres have degree $4$, and the
other two spheres have at most degree $3$'', and concludes from $2\cdot4 + 2\cdot3 = 14
< 16$. Neither assertion is justified, and the second is false: the graph $K_6$ minus a
perfect matching on four of its vertices has $13$ edges and degree sequence
$(5,5,4,4,4,4)$, satisfying every constraint derived up to that point. The preceding
step likewise yields only $\sum_{i\le4} \deg x_i \le 16$ against a requirement of
$\ge 16$, which is not a contradiction.

More fundamentally, the argument uses no geometric hypothesis at all: it manipulates
only degree sequences. No such argument can succeed, because a graph on six vertices
with $13$ edges exists. That $13$ contacts are impossible is a fact about $\R^3$, and
some input from the geometry of $\R^3$ is unavoidable --- in the present paper it is
Lemmas~\ref{lem:notri} and~\ref{lem:nosq}.

The arguments for $n = 7$ and $n = 8$ are case splits over degree sequences in which
every case is discharged by appeal to ``Lemmas 1--6''. Those lemmas carry the entire
geometric content of the paper and are stated without proof.

\bibliographystyle{plain}
\bibliography{refs}

\end{document}
